\chapter{Auditing Definitions and Declaration Bundles}\label{ch:definition-audit}

\href{cosine-primitive-graph.html}{The focused graph} retains one mathematical
conclusion and several declarations inside each broad node. A bundle is a
presentation choice, not a newly defined Lean object or an axiom. Its original
declaration names, source modules, and checked types remain visible.

\section{A namespace is not an input}
\code{ComputableAnalysis.CosinePrimitive.pi} is defined in
\code{CosinePrimitiveData.lean} as four times geometric arctangent at one.
The namespace \code{CosinePrimitive} does not refer to the primitive theorem.
In particular, the same independently defined pi is supposed to occur in the
endpoint statement. No cosine integral is used in its definition.

The audit traverses the full stored declaration closures for $A$, pi, $S$,
$C$, and the integral, and rejects any reference to the endpoint statement,
the endpoint program, or its proof routes. The pi closure also excludes
sine, cosine and the inverse provider. This is stronger than checking the
spelling of the definition alone. The raw inverse uses arctangent rectangles;
the provider also carries certificates mentioning earlier pi identities.
Those certificate dependencies are not numerical calls to pi.

\section{The exact first-quadrant algorithms}
Let $U_n$ be the literal output of \code{ClosedArctanInverse.raw (2*x)} at
stage $n$, and let $s(u)=2u/(1+u^2)$ and $c(u)=(1-u^2)/(1+u^2)$.
The readback theorems \code{TrigonometricReadback.sine\_stage} and
\code{TrigonometricReadback.cosine\_stage} are proved by reduction and state
\[
 S(x)_n=[s(U_n^-),s(U_n^+)],\qquad
 C(x)_n=[c(U_n^+),c(U_n^-)].
\]
The cosine endpoint order is reversed. These statements have the hypothesis
that $x$ is rational and $0\le x\le1/2$. The displayed wrappers return zero
outside this chart; they are not the global trigonometric extensions.

The literal integral stage reads only the fixed cosine sums and their
rational error radii. At $t=1/2$ its initial box is $[-1000,2001/2]$; the
separate endpoint program initially returns $[0,1/2]$. These concrete Lean
checks demonstrate that the programs are not definitionally identical.
Their eventual equality of represented values is the substantive theorem.

\section{Machine checks on an AI-curated view}
Each bundle's graph anchors must be among its displayed declarations. Arrows
contract actual declaration-reference paths, and each path is checked.
Anchors include the underlying sine and cosine constructors, not only wrapper
names: the Mathlib value bridges must visibly depend on those native
constructors. Definition readbacks and explanatory companions do not become
prerequisites of the measured proof terms merely by being displayed.

The \href{three-proofs/definition-audit.json}{definition and bundle audit}
records the checked closures and paths. It is not a claim that a contracted
graph is a minimal mathematical basis. Statement/type references and
proof/definition references are both part of the exported declaration graph;
a namespace prefix by itself is neither.
